\documentclass{article}
\usepackage[utf8]{inputenc}
\usepackage[T1]{fontenc}
\usepackage{geometry}
\geometry{
    a4paper,
    left=2.5cm,
    right=2.5cm,
    top=2.5cm,
    bottom=2.5cm
}
\usepackage{graphicx}
\usepackage{float}
\usepackage{hyperref}
\usepackage[english]{babel}
\usepackage{amsmath}
\usepackage{amssymb}
\usepackage{array}
\usepackage{booktabs}
\usepackage{caption}
\usepackage{subcaption}
\usepackage{enumitem}
\usepackage{fancyhdr}
\usepackage{titlesec}
\usepackage{url}
\usepackage{textcomp}
\usepackage{xcolor}
\usepackage{listings}
\usepackage{multirow}
\usepackage{longtable}
\usepackage{tabularx}
\usepackage{times}

% Set up page style
\pagestyle{fancy}
\fancyhf{}
\rfoot{\thepage}
\renewcommand{\headrulewidth}{0pt}

% Title formatting
\titleformat{\chapter}[display]{\normalfont\bfseries}{\chaptertitle\thechapter}{0pt}{\Huge}
\titlespacing*{\chapter}{0pt}{50pt}{40pt}

% Document info
\title{Template Document}
\author{Academic Support Agent}
\date{\today}

\begin{document}
\maketitle

\begin{center}
\thispagestyle{empty}
\onehalfspacing 
\begin{center} 
\includegraphics[scale=0.6]{logo.jpg} 
\end{center}

\vspace{2cm}
{\huge  lgebra Linear}\\
\vspace{2cm} % MUDAR ESTE VALOR PARA ALINHAR AP S MUDAR A DESCRI O DA TAREFA

\noindent\rule{8cm}{0.4pt}
\vspace{10mm}

{\Large Tarefa 5}\\
\vspace{5mm}
% DESCRI O DA TAREFA
{ UNIDADE 4 - TRANSFORMA ES LINEARES }\\
4.1 - Defini o e propriedades.\\
4.2 - Transforma es lineares no plano e no espa o.\\
4.3 - N cleo e imagem de uma transforma o linear.\\
4.4 - Matriz associada a uma transforma o linear.\\
4.5 - Opera es com transforma es lineares.\\
4.6 - Operadores lineares e operadores invers veis.\\
4.7 - Operadores ortogonais e sim tricos.\\
\vspace{10mm}

\noindent\rule{8cm}{0.4pt}\\
\vspace{2cm} % MUDAR ESTE VALOR PARA ALINHAR AP S MUDAR A DESCRI O DA TAREFA

{\large Prof. Jos  Ant nio Trindade Borges da Costa}\\
Tutora: Dion ia Migotto\\
Aluno: Tiago Cardoso Floriano\\
Matr cula: 202021740 %\hline
\end{center}

\newpage
\onehalfspacing 
\noindent
\textbf{ATIVIDADE 1}, 5) Dada a fun o $T:\mathbb{R}^2 \rightarrow \mathbb{R}^3$, definida por T(x, y) = (x+ky, x+k, y), verificar para quais dos seguintes valores de $k$, $T$   uma transforma o linear: \\
\\
a. $k = x$ \\
\\
{\textbf{SOLU O:}} \\
Tomando o exemplo 1 (p. 142) \cite{ines}, e considerando $\vec{u}=(x_1,y_1)$ e $\vec{v}=(x_2,y_2)$, teremos em $T(\vec{u}+\vec{v})$:\\
\begin{equation}
    \begin{split}
        T(\vec{u}+\vec{v})=T(x_1+x_2, y_1+y_2)=\\
        ((x_1+x_2)+(x_1+x_2)(y_1+y_2),(x_1+x_2)+(x_1+x_2), (y_1+y_2))=\\
        ((x_1+x_2)+(x_1 y_1)+(x_2 y_2),(2x_1+2x_2),(y_1+y_2))=\\
        ((2x_1 y_1+2x_2 y_2), (2x_1 +2x_2), (y_1+y_2))=\\
        ((2x_1 y_1, 2x_1, y_1) + (2x_2 y_2, 2x_2, y_2)=\\
        T(x_1,y_1)+T(x_2,y_2)=\\
        T(\vec{u}+\vec{v})
    \end{split}
\end{equation}



\newpage
\noindent 

\begin{thebibliography}{9}
%... items da bibliografia aqui

\bibitem{ines} FERREIRA, In s Farias,
	\emph{ lgebra Linear}, UAB-UFSM, 2010.
\bibitem{leon} LEON, Steven J.,
	\emph{ lgebra Linear com Aplica es}, 9th ed. Rio de Janeiro: LTC; 2018. ISBN: 978-85-216-3577-2.

\end{thebibliography}
\end{document}